\section{Discusión de resultados e implicancias para el diseño regulatorio}
\label{sec:implicancias_regulatorias}

Esta sección interpreta conjuntamente los resultados del análisis y examina sus implicancias para el diseño regulatorio de las metas de gestión. El objetivo no es derivar una arquitectura óptima ni formular una propuesta de implementación inmediata, sino identificar qué aspectos del diseño vigente resultan afectados por los mecanismos examinados, qué modificaciones encuentran mayor respaldo en el modelo y bajo qué condiciones podrían ser consideradas.

Los resultados no atribuyen una misma conclusión a todas las transiciones. La focalización de la medición financiera mejora la separación entre instrumentos, aunque su efecto sobre la ejecución contemporánea permanece condicionado; la eliminación del umbral agregado impide que el cumplimiento en otras dimensiones desactive consecuencias asociadas con incumplimientos ya configurados; la sustitución de la verificación financiera por indicadores físicos o funcionales evita que una reducción eficiente del gasto sea tratada, por sí sola, como menor cumplimiento de las intervenciones de \(R^{F}\); y la delimitación del ICG concentra su contenido informativo en los resultados que se pretende hacer visibles. Estas diferencias aconsejan discutir las implicancias de cada modificación por separado, antes que tratarlas como un paquete regulatorio indivisible.

La secuencia \(A_{0}\rightarrow A\rightarrow A'\rightarrow A''\rightarrow B\) cumple, por tanto, una función analítica. Permite identificar el mecanismo asociado con cada cambio y distinguir entre efectos sobre la ejecución, la eficiencia y la señal reputacional. La extensión \(B^{+}\) tiene un estatuto distinto: no corrige una superposición entre instrumentos, sino que incorpora un reconocimiento acotado del sobrecumplimiento y, por ello, requiere condiciones de elegibilidad y salvaguardas propias.


\subsection{Interpretación conjunta de los resultados}
\label{subsec:discusion_resultados}

La transición \(A_{0}\rightarrow A\) focaliza el avance financiero en las intervenciones de \(R^{F}\) y retira la relación de trabajo del ICG y de las consecuencias pecuniarias. Para las inversiones de \(R^{S}\) y \(R^{G}\), la focalización elimina su participación en el avance financiero general y, con ello, la posibilidad de que una reducción eficiente del costo deteriore la señal financiera o aumente la consecuencia asociada con dicho avance. Al mismo tiempo, la salida de la relación de trabajo elimina otra vía mediante la cual una misma inversión podía incidir en el índice y en la base de las consecuencias. Esta mayor separación no implica, sin embargo, un aumento incondicional del esfuerzo contemporáneo de ejecución: la transición modifica simultáneamente varios canales, por lo que el modelo no establece un orden general entre \(e_r^{A_0,*}\) y \(e_r^{A,*}\). Para las intervenciones de \(R^{S}\) y \(R^{G}\), en cambio, la focalización sí elimina los canales financieros que podían penalizar una reducción eficiente del costo, de modo que el incentivo de eficiencia no se debilita y puede aumentar cuando alguno de esos canales se encontraba activo.

La transición \(A\rightarrow A'\) modifica la regla de activación de las consecuencias pecuniarias. El ICG conserva la misma composición y las intervenciones de \(R^{F}\) continúan verificándose mediante el avance financiero focalizado, pero el umbral agregado deja de determinar si los incumplimientos ya configurados generan una consecuencia. Bajo \(A\), un valor \(ICG^{A}\geq\tau\) puede desactivar conjuntamente las consecuencias pese a que subsistan obligaciones incumplidas; desde \(A'\), cada componente conserva su consecuencia mientras permanezca configurado el incumplimiento correspondiente. Esta modificación puede fortalecer el incentivo de ejecución cuando el umbral habría mantenido inactivo el canal sancionador. Para \(R^{F}\), sin embargo, también puede reducir el incentivo de eficiencia al hacer efectiva una consecuencia financiera que bajo \(A\) podía encontrarse desactivada. La eliminación del umbral no sustituye el análisis de imputabilidad ni la necesidad de que la obligación y la conducta sancionable se encuentren previamente identificadas.

La transición \(A'\rightarrow A''\) sustituye la verificación financiera de \(R^{F}\) por indicadores específicos de cumplimiento físico o funcional de \(J^{F}\). El avance financiero focalizado deja de formar parte del ICG y de la base utilizada para determinar las consecuencias de estas intervenciones. En su lugar, el incumplimiento se verifica directamente mediante el indicador aplicable y la cuantificación se vincula con el alcance físico pendiente. Para las intervenciones de \(R^{S}\), la salida de \(f_F\) aumenta el peso de \(C_S\) dentro del ICG y fortalece su incentivo de ejecución. Para \(R^{F}\), el efecto sobre la ejecución es, en general, ambiguo, porque desaparecen tanto el retorno reputacional asociado con elevar el gasto como la reducción de la base financiera, mientras aparece el canal asociado con la base física. El efecto sobre la eficiencia es más definido: al retirarse los canales que penalizaban una reducción del gasto, \(A''\) no debilita y puede fortalecer el incentivo de eficiencia respecto de \(A'\).

La transición \(A''\rightarrow B\) completa la delimitación de la señal agregada. Al retirar \(C_G\) del ICG, el índice queda definido exclusivamente por \(C_S\), esto es, por el cumplimiento de los resultados comprendidos en \(J^{S}\). Para las inversiones de \(R^{S}\), esta modificación puede fortalecer el retorno reputacional del esfuerzo cuando la mayor concentración del índice en resultados incrementa suficientemente su legibilidad y valoración. Para las intervenciones de \(R^{G}\), en cambio, desaparece el retorno reputacional asociado con \(C_G\), aunque sus consecuencias específicas permanecen vigentes. Las intervenciones de \(R^{F}\) no modifican su tratamiento en esta transición, pues desde \(A''\) se encuentran fuera del ICG y sujetas a sus indicadores específicos de cumplimiento físico o funcional.

En conjunto, los resultados no establecen un ordenamiento global \(B\succ A_{0}\) para todos los márgenes de decisión. El aporte del modelo consiste, más bien, en identificar qué problema corrige cada transición y qué efectos puede producir sobre distintos incentivos. Esta lectura es especialmente importante para la discusión regulatoria: una modificación puede mejorar la congruencia entre objeto e instrumento sin aumentar necesariamente todos los esfuerzos observados, y un incentivo más intenso no es por sí mismo evidencia de una mejor arquitectura si proviene de una métrica que premia gasto redundante o penaliza eficiencia.


\subsection{Implicancias para la composición y función del ICG}
\label{subsec:implicancias_icg}

Si el ICG debe cumplir principalmente una función informativa y reputacional después de la desvinculación de los incrementos tarifarios base, los resultados sugieren que su composición debería partir del concepto que se pretende representar. Bajo la alternativa analizada, \(J^{S}\) comprende indicadores que miden directamente condiciones observables de cobertura y calidad del servicio, mientras que las inversiones, actividades, capacidades de gestión y demás medios utilizados para producirlas permanecen sujetos a instrumentos específicos.

Esta delimitación debe realizarse a partir del objeto directamente medido por cada indicador y no únicamente de su ubicación normativa dentro del SIIGEPSS. Los aspectos de cobertura y calidad pueden contener, junto con resultados propiamente dichos, variables que verifican ejecución física o funcional. Del mismo modo, indicadores de eficiencia empresarial como la relación de trabajo o el agua no facturada pueden ser relevantes para la supervisión y para la sostenibilidad de la prestación sin constituir, por ese solo hecho, resultados de \(J^{S}\). Su eventual exclusión del ICG no implicaría dejar de medirlos, sino separar su función de seguimiento de la señal agregada de desempeño.

La micromedición requiere una clasificación particularmente cuidadosa. Cuando el indicador verifica principalmente la instalación efectiva de medidores o el cumplimiento de un alcance físico identificable, su naturaleza se aproxima a una obligación de implementación y no a un resultado final de \(J^{S}\). Cuando, en cambio, se utilice un indicador que represente una condición operativa más amplia, su tratamiento deberá depender del objeto efectivamente medido. La pertenencia formal a una categoría normativa no debería sustituir esta evaluación funcional.

Las demás metas de gestión y dimensiones regulatorias pueden conservar reportes, metas, mecanismos de supervisión y, cuando corresponda, consecuencias autónomas. \(C_G\) resume en el modelo las metas de gestión distintas de los resultados del servicio, mientras que otros indicadores, como la relación de trabajo, el agua no facturada, los indicadores de sostenibilidad, solvencia, gobernanza, cumplimiento normativo o focalización del subsidio, pueden proporcionar información valiosa para fines específicos aun cuando no formen parte de \(C_S\). El cambio examinado consiste en evitar que todos ellos modifiquen automáticamente una medida destinada a resumir resultados de cobertura y calidad.

La utilidad de un ICG concentrado en resultados depende, además, de la difusión y legibilidad de la señal. Si el canal reputacional representado por \(d_I\) y \(\ell_I^{m}\) es débil, retirar componentes del índice no garantiza por sí mismo una respuesta gerencial relevante. Una eventual reforma tendría que acompañarse, por tanto, de una política de publicación que permita comparar resultados entre empresas y a lo largo del tiempo, mostrar los componentes del índice y distinguir con claridad el desempeño agregado de los demás indicadores regulatorios.

La evaluación individual y territorial debería mantenerse como complemento de la señal agregada. El ICG no sustituye la identificación del indicador específico incumplido ni la información necesaria para interpretar brechas entre localidades. La separación propuesta afecta la función del agregado, no la disponibilidad de información desagregada para fines de supervisión, transparencia o adopción de medidas correctivas.

La ponderación de los componentes también requiere una decisión explícita. Cuando el ICG se construye como un promedio simple, el número de metas asignadas a una dimensión determina implícitamente su incidencia agregada. Si se adopta una medida destinada a representar resultados de desempeño, la selección de ponderadores debería responder a un marco conceptual explícito y no únicamente a la cantidad de indicadores incluidos en cada estudio tarifario. Esta decisión es distinta de la valoración económica utilizada para determinar sanciones o reconocimientos y no tiene por qué reproducirla.


\subsection{Implicancias para la fiscalización y el régimen sancionador}
\label{subsec:implicancias_fiscalizacion_sanciones}

La exclusión de una inversión o de otra obligación de implementación del ICG no reduce su exigibilidad. Las intervenciones de \(R^{S}\), \(R^{G}\) y \(R^{F}\) continuarían sujetas a programación, seguimiento, fiscalización y, cuando exista una obligación autónoma, atribuible y expresamente tipificada, a medidas correctivas o consecuencias sancionadoras. La diferencia es que su tratamiento debe guardar correspondencia con el objeto de la obligación y no depender necesariamente de su incidencia dentro de una medida agregada de desempeño.

El Programa de Ejecución de Metas Físico-Financieras ofrece una base institucional para esta separación. Como instrumento de trazabilidad, puede vincular cada inversión o medida de mejora con su alcance, presupuesto, fuente de financiamiento, fases de ejecución, hitos, fecha prevista de entrada en operación, riesgos y medidas de contingencia. Esta información permitiría distinguir entre retrasos atribuibles, cambios de programación compatibles con el marco regulatorio y circunstancias externas que requieran un tratamiento específico.

Los resultados del modelo respaldan una distinción entre tres dimensiones que no deberían confundirse: alcance físico o funcional, costo y entrada efectiva en operación. El gasto realizado no acredita por sí mismo que el alcance haya sido completado; una inversión puede culminarse a un costo inferior al presupuesto; y una obra formalmente terminada puede no encontrarse todavía operativa. La fiscalización puede utilizar conjuntamente información financiera, física y funcional, pero la variable empleada para determinar el cumplimiento debería guardar correspondencia con la obligación efectivamente exigida.

El caso de EPS SEDA HUÁNUCO S.A., desarrollado en la discusión de los problemas de activación sancionadora de la Sección~3, ilustra la relevancia de esta distinción. La intervención examinada completó su alcance por debajo del presupuesto y el ahorro fue excluido de la base utilizada para determinar la multa, aunque el menor gasto continuó reduciendo el indicador financiero. En la representación del modelo, este problema corresponde al tratamiento de las intervenciones de \(R^{F}\): \(A\) y \(A'\) mantienen la verificación financiera focalizada, mientras que \(A''\) y \(B\) la sustituyen por indicadores específicos de cumplimiento físico o funcional. La implicancia no es que todo menor gasto constituya eficiencia, sino que un costo inferior no debería interpretarse por sí mismo como menor cumplimiento cuando se hayan preservado el alcance, la calidad, la capacidad, la vida útil y las demás condiciones exigibles.

La separación de instrumentos también obliga a revisar la relación entre el umbral agregado y las obligaciones específicas. Si una obligación permanece incumplida, el modelo muestra que condicionar su consecuencia al valor de un índice que agrega otras metas puede generar una zona en la que el mayor cumplimiento de unas dimensiones neutralice el incumplimiento de otras. La transición \(A\rightarrow A'\) elimina precisamente esta condición de activación. Esto justifica examinar si el umbral del ICG debe continuar operando como activador de consecuencias asociadas con obligaciones individualizables. La revisión no implica que toda brecha inferior al 100~\% deba sancionarse automáticamente: razones de materialidad, error de medición, costos de fiscalización o causas no atribuibles pueden justificar tolerancias o exclusiones, pero deberían tratarse mediante reglas congruentes con la obligación específica.

La imputabilidad conserva un papel central. Las causales previstas en el numeral 30.2 del Reglamento General de Fiscalización y Sanción permiten valorar circunstancias como caso fortuito, fuerza mayor, hecho determinante de tercero u otras situaciones previstas por la normativa. Bajo una arquitectura de obligaciones específicas, este análisis operaría después de verificar la existencia de la brecha y antes de atribuir responsabilidad. De este modo, separar la consecuencia del ICG no equivale a convertir automáticamente cualquier incumplimiento físico en responsabilidad administrativa.

La base económica de la sanción debería guardar correspondencia con la conducta tipificada y evitar valorizaciones repetidas. Cuando una misma inversión contribuya a varios resultados o se encuentre relacionada con varias metas, su valor no debería contabilizarse íntegramente más de una vez dentro de la consecuencia económica. El Anexo~\ref{app:calculo_sancionador} desarrolla esta correspondencia para las inversiones no ejecutadas. Esta cuestión es distinta del \textit{non bis in idem}: la duplicación de una base económica puede ocurrir incluso dentro de una sola multa, mientras que el principio de \textit{non bis in idem} controla la imposición de más de una sanción por una misma conducta cuando concurren las identidades exigidas por la normativa.

Los incumplimientos de resultados de \(J^{S}\) requieren una lógica propia. Una meta de desempeño puede generar la consecuencia específica prevista para su incumplimiento sin que esta dependa del valor agregado del ICG. Además, una misma dimensión puede encontrarse sujeta a un estándar normativo autónomo de continuidad, presión, calidad, acceso u otra condición de la prestación. En ese caso, la consecuencia debería responder a la afectación producida, al beneficio ilícito o a otra base económica congruente con la obligación incumplida, y no reutilizar automáticamente el costo de una inversión ya considerado dentro del control del programa de inversiones.

La misma lógica se extiende a las demás metas de gestión y a los indicadores operativos que permanecen fuera de \(C_S\). La relación de trabajo, el agua no facturada u otros indicadores pueden conservar exigibilidad cuando exista una obligación autónoma, pero su eventual fiscalización debe corresponder a la conducta que representan. No resulta necesario reproducir para ellos la metodología de alcance físico aplicable a una intervención discreta cuando su producción depende de decisiones continuas de operación, gestión comercial, organización e inversión.

Las dificultades descritas pueden reducirse mediante una tipificación vinculada directamente con el incumplimiento atribuible del programa de inversiones. Bajo este enfoque, la conducta sancionable sería la falta de ejecución oportuna de una inversión previamente identificada, programada y financiada.

La obligación incumplida correspondería a la inversión \(r\); el monto relevante sería \(U_r\), construido conforme al procedimiento de determinación del monto dejado de invertir desarrollado en el Anexo~\ref{app:calculo_sancionador}; y la consecuencia económica se calcularía a partir de su costo postergado \(CP_r\), definido en la ecuación~\eqref{eq:anexo_costo_postergado}. De este modo, cada inversión sería valorada una sola vez, con independencia del número de metas o indicadores a los que contribuya.

Esta formulación establece una correspondencia directa entre la conducta tipificada, la base económica y la sanción. Asimismo, elimina la necesidad de reconstruir índices individuales hipotéticos para distribuir la brecha de un ICG agregado cuando la obligación incumplida puede identificarse directamente en el programa de inversiones.

La tipificación basada en inversiones no impide mantener estándares ni consecuencias específicas frente a incumplimientos autónomos de los resultados del servicio. Si una reducción de la continuidad, la presión, la calidad del agua u otro indicador genera una afectación independiente a los usuarios, dicha conducta puede evaluarse mediante medidas correctivas o tipificaciones específicas de calidad. En ese supuesto, la consecuencia regulatoria debe guardar relación con la afectación producida y no reutilizar como base económica la inversión ya considerada en el incumplimiento del programa.

En síntesis, la metodología basada en índices parte del incumplimiento de una meta o de un indicador agregado y posteriormente identifica las inversiones vinculadas con dicho incumplimiento. Este procedimiento resulta relativamente directo cuando existe una relación clara entre una meta y una inversión, pero se vuelve más complejo cuando la infracción se configura mediante el ICG o cuando existen inversiones compartidas. Una tipificación centrada en el programa de inversiones permitiría identificar directamente cada obligación incumplida, valorar una sola vez su postergación y reducir la necesidad de reconstrucciones contrafactuales y asignaciones convencionales.


\subsection{Condiciones para una eventual implementación}
\label{subsec:condiciones_eventual_implementacion}

Las implicancias anteriores no determinan por sí mismas una reforma inmediata. Una eventual implementación requeriría, en primer lugar, revisar el sistema de indicadores para clasificar cada variable según el objeto predominantemente medido: resultado de \(J^{S}\), intervención vinculada con \(R^{S}\), \(R^{G}\) o \(R^{F}\), indicador de eficiencia empresarial u otro propósito regulatorio especializado. Esta revisión permitiría evitar que la ubicación normativa de un indicador determine automáticamente su tratamiento dentro del ICG.

En segundo lugar, las obligaciones de implementación tendrían que encontrarse suficientemente individualizadas. Para fiscalizar separadamente una inversión, actividad o medida de mejora es necesario identificar su alcance exigible, plazo, hitos, medios de verificación, condiciones de entrada en operación, presupuesto o referencia económica y tratamiento de causas no atribuibles. La separación del ICG solo es sostenible si el instrumento alternativo de control conserva trazabilidad suficiente.

En tercer lugar, sería necesario asegurar la consistencia entre los estudios tarifarios vigentes, las reglas de fiscalización y las tipificaciones. La transición plantea una dificultad especial porque pueden coexistir estudios aprobados con metas de avance financiero de alcance amplio y otros sujetos a la focalización introducida posteriormente. Cualquier modificación adicional debería definir expresamente el tratamiento de las obligaciones ya aprobadas y evitar alterar retroactivamente las condiciones bajo las cuales se fijaron tarifas, metas y programas de inversión.

En cuarto lugar, la utilidad de un ICG predominantemente reputacional depende de la calidad de la información y de su difusión. La separación de componentes solo mejora la señal si el regulador publica resultados de manera oportuna, comparable y comprensible, y si los actores relevantes pueden distinguir entre desempeño, ejecución de inversiones y otras dimensiones de gestión. Sin esta condición, la mejora en la composición formal del índice podría no traducirse en un incentivo reputacional material.

En quinto lugar, la reforma del régimen sancionador requeriría información que permita relacionar cada obligación con el alcance pendiente, el presupuesto reconocido, el estado de ejecución, la entrada en operación y las causas del retraso. Sin esa trazabilidad, sustituir una regla agregada por consecuencias individualizadas podría aumentar la discrecionalidad o los costos de fiscalización. La separación de instrumentos exige, por tanto, capacidad administrativa suficiente para sostener el mayor grado de identificación que propone.

La secuencia \(A_{0}\rightarrow A\rightarrow A'\rightarrow A''\rightarrow B\) no debe interpretarse como un calendario obligatorio de reforma. \(A_{0}\rightarrow A\) reconstruye una modificación normativa que focaliza el avance financiero y, en la representación analítica, retira la relación de trabajo del ICG y de las consecuencias pecuniarias; las transiciones posteriores constituyen comparaciones de diseño. La implementación podría ser gradual y no necesariamente reproducir exactamente el orden del modelo, pero cualquier combinación debería preservar la coherencia entre medición, fiscalización y sanción. En particular, excluir obligaciones del ICG sin asegurar previamente su seguimiento y exigibilidad separada podría debilitar el incentivo de ejecución.


\subsection{Tratamiento del sobrecumplimiento}
\label{subsec:implicancias_bonificacion}

La extensión \(B^{+}\) debe interpretarse separadamente de las modificaciones anteriores. Mientras \(A_{0}\rightarrow B\) analiza la congruencia entre objetos e instrumentos, \(B^{+}\) incorpora un incentivo adicional destinado a reconocer resultados verificables por encima de la meta. Su eventual utilización requiere, por tanto, una evaluación específica y no constituye una consecuencia necesaria de delimitar el ICG.

El cumplimiento ordinario puede continuar truncándose en 100~\%, de modo que el componente base \(C_S\) no permita que el sobrecumplimiento de una meta eleve el cumplimiento ordinario de otra. \(B^{+}\) identifica separadamente el desempeño adicional y lo incorpora de manera acotada al ICG extendido, además de asociarle, cuando corresponda, una bonificación predeterminada. Como las consecuencias específicas dejan de depender del nivel agregado del ICG desde \(A'\), el sobrecumplimiento reconocido no desactiva incumplimientos ni consecuencias correspondientes a otras obligaciones.

La selección de los resultados elegibles tendría que considerar, entre otros elementos, la relevancia para los usuarios, la controlabilidad, la verificabilidad, la sostenibilidad de la mejora y la posibilidad de atribuirla razonablemente a las decisiones de la empresa. También deberían definirse ex ante el rango máximo de sobrecumplimiento reconocido, la valoración aplicable por unidad adicional y las condiciones que excluyen el reconocimiento cuando el resultado responde principalmente a factores externos.

El modelo abstrae de restricciones vinculantes de capacidad gerencial y supone costos separables entre intervenciones para aislar el efecto del reconocimiento del sobrecumplimiento. Bajo esos supuestos, \(B^{+}\) no desplaza mecánicamente esfuerzo desde otras actividades y puede mantener un retorno marginal positivo una vez alcanzada la meta, dentro del rango reconocido. En la práctica, sin embargo, una capacidad gerencial o unos recursos limitados podrían generar reasignaciones desde obligaciones todavía pendientes. Este riesgo queda fuera del alcance formal del modelo y debería considerarse al diseñar las reglas de elegibilidad y los límites del mecanismo.

El reconocimiento puede ser reputacional, regulatorio o económico. Una modalidad reputacional resulta consistente con la función informativa del ICG y puede implementarse mediante publicaciones, comparaciones o reconocimientos diferenciados. Una modalidad económica exigiría una evaluación adicional de financiamiento, incidencia tarifaria, distribución de beneficios entre empresa y usuarios y límites por indicador o periodo. En ningún caso el reconocimiento debería restablecer automáticamente la condicionalidad general de los incrementos tarifarios base ni operar como mecanismo de compensación de incumplimientos en otras dimensiones.

La gradualidad es particularmente relevante para \(B^{+}\). Antes de introducir una bonificación económica generalizada, podría evaluarse un conjunto acotado de indicadores con información suficientemente robusta, reglas de elegibilidad predeterminadas y mecanismos para verificar persistencia y atribución. La continuidad o ampliación del mecanismo debería depender de que la evidencia obtenida permita distinguir mejoras adicionales de variaciones transitorias o de resultados producidos principalmente por factores externos.

En síntesis, las implicancias del análisis apuntan menos a la adopción de un paquete cerrado que a una separación funcional entre objetos regulatorios. El ICG puede concentrarse en los resultados de \(J^{S}\); las intervenciones de \(R^{S}\), \(R^{G}\) y \(R^{F}\) pueden permanecer sujetas a programación, fiscalización y consecuencias específicas conforme a la naturaleza de cada obligación; los demás indicadores regulatorios pueden conservar mecanismos especializados de seguimiento; y el sobrecumplimiento puede tratarse, si se justifica, mediante un componente adicional y acotado. La conveniencia de cada modificación depende de las condiciones de información, controlabilidad, atribución y capacidad institucional desarrolladas en esta sección.


\begin{figure*}[t]
\centering

\begin{tikzpicture}[
font=\small,
cadena/.style={
    draw,
    rounded corners=2pt,
    align=center,
    minimum height=2.25cm,
    text width=0.155\textwidth,
    inner sep=5pt
},
instrumento/.style={
    draw,
    rounded corners=2pt,
    align=center,
    minimum height=1.65cm,
    text width=0.155\textwidth,
    inner sep=5pt
},
flecha/.style={
    -{Latex},
    thick
},
control/.style={
    -{Latex},
    densely dashed,
    thick
},
frontera/.style={
    draw,
    dashed,
    rounded corners=4pt,
    inner sep=5pt
}
]

% Cadena principal
\node[cadena] (recursos) {
\textbf{Recursos}\\[2mm]
Recursos tarifarios y fuentes de financiamiento
};

\node[cadena, right=5mm of recursos] (implementacion) {
\textbf{Obligaciones de implementación}\\[2mm]
Inversiones, actividades y medidas de mejora\\[1mm]
\(\left(R^{S}\mathbin{\dot\cup}R^{G}\mathbin{\dot\cup}R^{F}\right)\)
};

\node[cadena, right=5mm of implementacion] (operacion) {
\textbf{Alcance y operación}\\[2mm]
Alcance físico o funcional completado y entrada efectiva en operación
};

\node[cadena, right=5mm of operacion] (resultados) {
\textbf{Resultados de desempeño}\\[2mm]
Cobertura y calidad del servicio\\[1mm]
\(\left(J^{S}\right)\)
};

\node[cadena, right=5mm of resultados] (sobrecumplimiento) {
\textbf{Desempeño adicional}\\[2mm]
Sobrecumplimiento verificable y acotado
};

% Flechas de la cadena
\draw[flecha]
(recursos.east) -- (implementacion.west);

\draw[flecha]
(implementacion.east) -- (operacion.west);

\draw[flecha]
(operacion.east) --
node[
    above,
    align=center,
    font=\scriptsize
] {contribución productiva,\\no equivalencia automática}
(resultados.west);

\draw[flecha]
(resultados.east) -- (sobrecumplimiento.west);

% Frontera del ICG base
\node[
    frontera,
    fit=(resultados),
    label={
        [font=\scriptsize]
        above:\textbf{ICG base de desempeño}
    }
] (icgbox) {};

% Instrumentos regulatorios
\node[instrumento, below=1.65cm of recursos] (seguimiento) {
\textbf{Programación y seguimiento}\\[2mm]
Programación físico-financiera, cronogramas, hitos y alertas
};

\node[instrumento, below=1.65cm of implementacion] (fiscalizacion) {
\textbf{Fiscalización de obligaciones}\\[2mm]
Verificación del cumplimiento y análisis de atribuibilidad
};

\node[instrumento, below=1.65cm of operacion] (verificacion) {
\textbf{Verificación física y funcional}\\[2mm]
Culminación, calidad técnica y entrada efectiva en operación
};

\node[instrumento, below=1.65cm of resultados] (evaluacion) {
\textbf{Evaluación del desempeño}\\[2mm]
ICG, componentes y publicación de resultados
};

\node[instrumento, below=1.65cm of sobrecumplimiento] (bonificacion) {
\textbf{Reconocimiento adicional}\\[2mm]
Componente acotado y bonificación bajo \(B^{+}\)
};

% Flechas verticales
\draw[control]
(seguimiento.north) -- (recursos.south);

\draw[control]
(fiscalizacion.north) -- (implementacion.south);

\draw[control]
(verificacion.north) -- (operacion.south);

\draw[control]
(evaluacion.north) -- (resultados.south);

\draw[control]
(bonificacion.north) -- (sobrecumplimiento.south);

% Etiquetas laterales
\node[
    rotate=90,
    font=\footnotesize\bfseries,
    left=7mm of recursos
] {Cadena de producción};

\node[
    rotate=90,
    font=\footnotesize\bfseries,
    left=7mm of seguimiento
] {Instrumento regulatorio};

\end{tikzpicture}

\caption{Cadena de resultados y separación funcional de instrumentos regulatorios}
\label{fig:cadena_resultados_instrumentos}

\begin{minipage}{0.97\textwidth}
\footnotesize
\textit{Nota:} Las flechas horizontales representan la cadena mediante la cual los recursos y las obligaciones de implementación contribuyen a producir resultados. Esta relación no implica equivalencia entre la ejecución de los medios y la obtención del desempeño. El ICG base se concentra en los resultados comprendidos en \(J^{S}\), mientras que las intervenciones de \(R^{S}\), \(R^{G}\) y \(R^{F}\) permanecen sujetas a programación, seguimiento, fiscalización y, cuando corresponda, consecuencias específicas. Otras dimensiones regulatorias que no se representan en la cadena principal conservan mecanismos especializados de seguimiento y supervisión. Bajo \(B^{+}\), el sobrecumplimiento se identifica mediante un componente adicional y acotado del ICG extendido y puede recibir una bonificación predeterminada, sin modificar el cumplimiento base ni desactivar consecuencias correspondientes a otras obligaciones. \textit{Fuente:} Elaboración propia sobre la base de \citet{oecd2023glossary}, \citet{behn2003measure} y la experiencia regulatoria revisada en el Anexo~\ref{anexo:experiencia_internacional}.
\end{minipage}

\end{figure*}